\section{Conclusion and Future Work}

In this work, we analyze the efficacy of using multi-agent debate at inference time to guard LLMs against adversarial attacks. We implement multi-agent debate settings between a variety of different models with different intentions. We then prompt models equipped with multi-agent feedback with previously successful adversarial attacks on models, finding that multi-agent debate does appear to have some effect on lowering response toxicity when compared to baselines such as Self-Refine. However, we also find that an LLM agent that outputs toxic content may negatively influence other LLM agents it is interacting with in a debate context, although this effect is not as strong as the positive effect from multi-agent debate with non-poisoned agents. Additionally, this effect is also sensitive to the specific model and intent used.

While multi-agent debate shows some promise as an inference-time approach to safeguarding against adversarial prompts, our analysis shows that it is far from perfect. One weakness of such an approach is that querying larger models multiple times in a debate context can be very resource-intensive in both model cost and latency. However, if a model is not generally capable enough, it may lack the conversational ability to benefit from multi-agent debate. Future work could consider whether debate between models from different providers could help improve the adversarial robustness of each individual model, as well as what debate dynamics look like over longer interactions with more agents. This requires designing more sophisticated "debate" or "discussion" frameworks with realistic feedback and update stages, and potentially fine-tuning models on specific intention. Additionally, research into learning when longer debate is needed to steer a model away from toxic output could make this a more efficient and effective model guardrail.